\documentclass[11pt,a4paper]{article}
\usepackage[margin=25mm]{geometry}
\usepackage{amsmath,amssymb,amsthm,booktabs,longtable,array,hyperref}
\setlength{\emergencystretch}{3em}
\newtheorem{theorem}{Theorem}
\newtheorem{proposition}[theorem]{Proposition}
\newcommand{\gen}[1]{\langle #1\rangle}
\newcommand{\Aut}{\operatorname{Aut}}
\newcommand{\im}{\operatorname{im}}
\newcommand{\R}{\mathbb R}
\title{Chambers, universal constructions, and realization cones}
\author{Human Compression project}\date{}
\begin{document}\maketitle
\begin{abstract}
We organize the mathematical content of Schulte and Ivi\'c Weiss's
2006 survey \emph{Problems on Polytopes, Their Groups, and Realizations}.
Three common arguments treat distinguished generators and closed flag
walks, universal amalgamation with regular local types, and finite
equivariant realizations. Two oriented amalgamations are verified by
their complete presentations. Previously established group results
answer the sporadic existence question and the regular alternating
rank question within their precise scopes. The source's remaining
classification inputs and questions are recorded separately. In
particular, its two aggregate counts for locally projective polytopes
are reported claims, not reproduced censuses. All four component
manuscripts form part of the proof; this synthesis does not replace
their proofs by numerical summaries.
\end{abstract}

\section{The source and the mathematical interfaces}

The fixed source is arXiv:math/0608397v1, dated 15 August 2006,
corresponding to the paper in \emph{Periodica Mathematica Hungarica}
53 (2006), 231--255. We use its twenty-five physical pages and its
complete TeX source. A later compilation date on the acquired PDF
does not constitute a later mathematical edition. This is a survey
with thirty-four numbered problems, rather than a paper with a
single new classification theorem. Fourteen additional records
separate embedded conjectures, restrictions, and unnumbered research
directions. They are not fourteen additional independent conjectures.

An abstract polytope of rank $n$ is a ranked poset with unique least
and greatest faces, the diamond condition, and strong flag
connectedness. The diamond condition says that each section of
rank two has exactly two intermediate faces. Strong flag
connectedness requires a path between two flags that retains
their common faces. These conditions concern actual flags and
sections; they cannot be replaced by connectedness of an
unrelated incidence graph.

For a regular polytope, its distinguished involutions
$\rho_0,\ldots,\rho_{n-1}$ satisfy the string commutations and
\[
 \gen{\rho_i:i\in I}\cap\gen{\rho_j:j\in J}
       =\gen{\rho_k:k\in I\cap J}.
\]
For a chiral polytope, adjacent flags are in the two different
automorphism orbits. Its distinguished rotations
$\sigma_1,\ldots,\sigma_{n-1}$ satisfy the string rotation relations
and their rotation intersection condition. We retain the standard
marked reconstruction theorems, attributed in the survey to
McMullen--Schulte and Schulte--Weiss. Relations without the
intersection condition do not supply a polytope. The geometric
reconstruction is an explicit antecedent of the arguments below.

The complete general proof for highly symmetric hypertopes supplies
the corresponding stabilizer argument. Paths within a residue identify
its pointwise stabilizer with the subgroup on its remaining types.
In the chiral case even paths give the pair-difference subgroup.
Polytope flags are chambers of their face geometry, so the argument
applies here. Conversely, a general C-group or C$^+$-group needs the
geometric hypotheses stated in that proof. We use the marked string
reconstruction where appropriate, not an unrestricted converse for
arbitrary diagrams.

\section{Words, orientations, and universal objects}

\begin{theorem}\label{words}
For a regular or chiral polytope of rank $n\ge3$, put
$\alpha_0=1$ and $\alpha_i=\sigma_1\cdots\sigma_i$ for
$1\le i\le n-1$. Simultaneous inversion of the $\alpha_i$ is
equivalent to the usual reflection substitution
\[
 \sigma_1\longmapsto\sigma_1^{-1},\qquad
 \sigma_2\longmapsto\sigma_1^2\sigma_2,\qquad
 \sigma_i\longmapsto\sigma_i\quad(i\ge3).
\]
For a chiral polyhedron, let $s_1,s_2$ denote its two distinguished
rotations with the traversal convention in the component proof.
The length of a $j$-hole is
\[
                  \bigl|s_1s_2^{1-j}\bigr|,
\]
and the length of a $j$-zigzag is
\[
             2\bigl|s_1s_2^{-j}s_1^{-1}s_2^j\bigr|.
\]
Here $|g|$ denotes group order, and the traversal state includes
the selected side of an edge.
\end{theorem}
\begin{proof}
The complete three-page component \emph{Rotation conventions and
closed flag walks} proves these statements by word cancellation
and actual flag actions. The relevant identities are
$\sigma_i=\alpha_{i-1}^{-1}\alpha_i$ and $\alpha_i^2=1$ for
$i\ge2$. If $r_i$ is the operation of changing the type-$i$ face
of a flag, the traversal words are
\[
 h_j=r_0r_1(r_2r_1)^{j-1},\qquad z_j=r_0(r_1r_2)^j.
\]
These $r_i$ are flag permutations, not generally polytope
automorphisms. An automorphism commutes with every adjacency
operation. Consequently its iterates agree with successive
repetitions of the chosen word on a base flag whenever one
traversal returns to its automorphism orbit. The free action
on that orbit identifies the traversal period with its group
order. In a chiral action $h_j$ preserves the orbit whereas
$z_j$ exchanges the two orbits. The component computes
$h_j=s_1s_2^{1-j}$ and
$z_j^2=s_1s_2^{-j}s_1^{-1}s_2^j$ on the base flag. A return
under $z_j$ must therefore take an even number of steps,
which gives the second formula. The same component gives
the regular formulas and all convention changes explicitly.
\end{proof}
The formulas determine periods of a specified polyhedron.
They do not determine which period lists determine a polyhedron,
or when a universal quotient prescribed by those lists is finite.
Those are the separate assertions in Problem~7.

\begin{theorem}\label{universal}
Let $K,L$ be regular rank-$n$ polytopes, $n\ge2$, with
compatible distinguished common section. If any rank-$(n+1)$
polytope has facets $K$ and vertex figures $L$, there is a
regular universal polytope with these local types.
\end{theorem}
\begin{proof}
The full two-page component \emph{Universal regular amalgamation
from a witness} gives the argument, including the group-action
justifications. Start with involutions on the rank-$(n+1)$
string, imposing remote commutation. Add the complete marked
relations of $K$ on the first $n$ generators and of $L$ on
the last $n$. In a given witness the adjacency permutations
satisfy these relations at every flag: the local constituents
are regular, so their marked relations hold throughout each
local flag set. The presented group therefore acts on all
witness flags.

The two end subgroups identify with the exact local groups.
Surjections in one direction come from the local presentations;
their actions on the corresponding local flags give the reverse
order-free identifications. An element in both end subgroups
preserves the base facet and vertex. Strong flag connectedness
of their common section supplies a middle-subgroup element with
the same effect on the base flag. The free action of either
end group makes these two elements equal. Thus the end
intersection is exactly the middle subgroup. The retained
string end criterion and reconstruction theorem give a regular
polytope. Repeating the action on every other witness supplies
the covering map from this one, and proves universality.
\end{proof}
The witness assumption is essential. The source already gives
a pair with no regular amalgamation: facets $\{4,4\}_{(3,0)}$
and hemicube vertex figures. That negative example is retained
with its attribution to \emph{Abstract Regular Polytopes},
p.~367. Theorem~\ref{universal} supplies neither existence
for arbitrary compatible pairs nor a finiteness criterion.
The analogous theorem for oriented chiral local constituents
is a separate retained Schulte--Weiss input. It requires the
chosen enantiomorphic forms and directly regular constituents
where a constituent is regular.

\begin{proposition}\label{oriented}
The two oriented universal amalgamations with local pairs
\[
 (\{4,4\}_{(1,3)},\{4,4\}_{(1,3)})\quad\hbox{and}\quad
 (\{4,4\}_{(3,1)},\{4,4\}_{(1,3)})
\]
are chiral and have full automorphism groups of orders $2000$
and $960$, respectively.
\end{proposition}
\begin{proof}
The complete two-page component \emph{Two oriented square-torus
amalgamations} verifies the full presentations
\begin{align*}
 \gen{u,v,w\mid{}&u^4=v^4=w^4=(uv)^2=(vw)^2=(uvw)^2=1,\\
 &(u^{-1}v)^s(uv^{-1})^t=(v^{-1}w)(vw^{-1})^3=1}
\end{align*}
at $(s,t)=(1,3),(3,1)$. Exact isomorphisms from the finitely
presented groups to permutation groups give the stated orders.
The two end groups have order $40$ and their intersection is
$\gen v$ of order $4$; both adjacent cyclic intersections
are trivial. Separate local presentations also have order $40$,
so the local types do not collapse. Both toroidal relators
fail after the reflection substitution of Theorem~\ref{words}.
The retained rotation reconstruction therefore gives chirality.
The presentation imposes exactly the local relations, so its
universal property is the retained oriented amalgamation
correspondence. Independent literal permutation closure checks
the finite orders and intersections; the upper order of the
presentation itself remains the explicit GAP isomorphism
calculation. These established examples are attributed by
the survey to Hubard--Weiss (2005), p.~132.
\end{proof}

\section{Finite equivariant realization theory}

Let a finite group $G$ act on a nonempty finite set $X$. A realization is
a map $f:X\to V$ to a real orthogonal $G$-module with
$f(gx)=gf(x)$. We allow repeated vertices until faithfulness
is imposed, translate to centroid zero, and take the actual
linear span of the image. Write $\mathbf1$ for the all-ones
vector and $J=I-|X|^{-1}\mathbf1\mathbf1^T$.

\begin{theorem}\label{cone}
Congruence classes of centered equivariant realizations are
identified with
\[
 \mathcal C_G(X)=\{B=B^T\succeq0:B\mathbf1=0,\ gB=Bg
                         \text{ for every }g\in G\},
\]
where $g$ acts by its permutation matrix on $\R^X$.
The dimension of the realized span is $\operatorname{rank}B$.
Squared distances and Gram matrices are mutually recoverable by
\[
 D_{xy}=B_{xx}+B_{yy}-2B_{xy},\qquad B=-\tfrac12 JDJ.
\]
Blending realizations is addition in this cone. A nonzero
$B$ spans an extreme ray precisely when $\im B$ is an
irreducible real $G$-module.
\end{theorem}
\begin{proof}
The complete four-page component \emph{Finite realization cones}
proves the Gram correspondence, the inverse on hollow symmetric
distance matrices, and the extreme-ray statement. The realization
of $B$ is $f_x=B^{1/2}e_x$ in its actual image. Orthogonal
direct sums give addition of Gram matrices; one then restricts
to the span rather than assuming that the entire direct sum
is spanned.

For clarity, the point where irreducibility enters is as follows.
If $B=C+D$ with $C,D\succeq0$, then $\ker B\subseteq\ker C$,
so $C$ restricts to $\im B$. An invariant self-adjoint
operator on an irreducible real module is scalar: each of
its real eigenspaces is invariant. This applies to both $B$
and $C$, and makes $C$ a scalar multiple of $B$. Conversely,
if $\im B$ is reducible, choose an orthogonal invariant
projection $Q$ onto a nonzero proper summand. With
$A=B^{1/2}$ the decomposition
\[
                    B=AQA+A(I-Q)A
\]
has two nonzero invariant positive semidefinite terms.
Their ranks preclude either being a multiple of $B$.
This argument does not assume that $Q$ commutes with $B$.
\end{proof}

For a transitive vertex set $X=G/H$, all such orbit maps have
the form $gH\mapsto gv$ with $v\in V^H$. Distinct abstract
vertices have distinct images exactly when
$\operatorname{Stab}_G(v)=H$. Faithfulness of the representation
alone does not imply this equality. For a polytope one must also
separate distinct faces of each rank under the recursive face
construction. The component gives the precise recursive test,
and identifies the full geometric symmetry group as the Gram
stabilizer in the abstract automorphism group, acting on the
actual affine span.

The finite-group assumption is used, in particular, in averaging
and in the orthogonal representation setting. The source's
infinite Euclidean Wythoff constructions, discreteness conditions,
and geometric classifications are retained separately. Theorem
\ref{cone} gives neither all individual toroid cones nor the
answer about quaternionic irreducible characters in Problem~23.
It also does not decide whether an abstract regular polytope
has a faithful chiral geometric realization. The realization
theory retains the attribution to McMullen and to
McMullen--Monson recorded in the source; a new proof of its
common linear algebra is not a claim of priority for that theory.

\section{Group consequences already available}

The acquired simple-group theorem of Leemans--Liebeck (2016)
constructs a chiral polyhedron for every nonabelian finite
simple group outside $A_7$, $\mathrm{PSL}_2(q)$,
$\mathrm{PSL}_3(q)$, and $\mathrm{PSU}_3(q)$. We use its
positive assertion only. All twenty-six sporadic groups lie
in its scope, including the Monster, so it answers the
existence request in Problem~31. This is a previously
established answer, with its original attribution, not a
new solution obtained here or an enumeration of all possible
polytopes with those groups.

The acquired alternating-group existence and upper-rank
theorems, together with their complete finite boundary
censuses, give the following regular ranks for $n\ge3$:
\[
\begin{array}{c|l}
 n&\text{ranks of string C-markings generating }A_n\\\hline
 3,4,6,7,8&\text{none}\\
 5&3\\
 9&3,4\\
 10&3,4,5\\
 11&3,6\\
 n\ge12&3,4,\ldots,\lfloor(n-1)/2\rfloor.
\end{array}
\]
The exceptional gaps at $n=11$ are retained from the
complete finite census; they are not filled by interpolation.
The lower existence theorem and upper bound have separate
source attributions in their acquired manuscripts. The
simple-group theorem also supplies chiral rank-three
examples for $A_n$, $n\ge8$. These statements do not
classify all chiral ranks or all other polytopes requested
in Problem~32.

For Problem~30, the acquired common quotient for a
nondegenerate finite string C-marking of a $2$-group in
rank $d\ge2$ has order $2^{2d-1}$ and nilpotency class at
most two. The eight acquired families with large adjacent
generator orders have their stated parameter hypotheses.
Neither result classifies every $2$-group, every chiral
marking, or the groups of order $2^kp$ in the source question.

\begin{proposition}\label{neighborly}
The full automorphism group of a finite neighborly regular
or chiral polytope of positive rank acts $2$-transitively on its vertices.
\end{proposition}
\begin{proof}
First suppose the rank is at least three. Given an ordered
pair $(v,w)$ of distinct vertices, neighborliness supplies
an edge $e$ with these endpoints. Extend the incident
vertex-edge pair $(v,e)$ to a flag. In the regular case
all flags lie in one orbit. In the chiral case, if the flag
is in the unwanted orbit, change only its type-$2$ face.
This preserves $v$ and $e$ and moves the flag into the
chosen orbit. Do the same for any second ordered vertex
pair $(v',w')$. An automorphism taking the resulting first
flag to the second takes $v$ to $v'$ and $e$ to $e'$.
An edge has exactly two vertices by the diamond condition,
so it also takes $w$ to $w'$. Thus all ordered pairs of
distinct vertices are in one orbit. In rank one the two
vertices are interchanged. In rank two a neighborly
polygon has two or three vertices, and its dihedral
action is $2$-transitive on them. Chiral polytopes have
rank at least three by definition.
\end{proof}
This proves the necessary condition stated after Problem~33.
It is not a classification or a converse, and is not
claimed as a new theorem in the subject.

\section{Retained classifications and precise limits}

The survey collects results from many different sources.
The accompanying register has fifty-one clauses in sixteen
contexts, with source pages, literal anchors, citation keys,
and each acquired proof or retained premise. Its sixty-six
bibliography entries preserve the complete source list.
The following table describes the division of work. A
retained classification is an attributed input, not a
classification reproved by the arguments above.

\input{context_table.tex}

There are two numerical limits of particular importance.
The source reports seventeen universal locally projective
regular $4$-polytopes, eight with a $2$ in their Schl\"afli
symbol, and $441$ quotients including four additional
nonuniversal regular polytopes. These are attributed to
Hartley's then-preprint classification. We have not
reproduced the global lists or reconciled the complete
preprint and published editions. They therefore remain
reported aggregates, not verified conclusions of this paper.

One constituent is acquired in full: the universal
$\{\{5,3\},\{3,5\}_5\}$ has group
$J_1\times\mathrm{PSL}_2(19)$, and its accepted corrected
quotient census has $137$ entries under its specified
equivalence. It would be invalid to subtract a correction
for this family from $441$ without the original global
partition, overlap conventions, and edition correspondence.
No corrected global total is asserted.

The source's report of map enumerations through orientable
genus $15$ and nonorientable genus $30$ includes a chiral
claim whose cited paper is titled \emph{Determination of
all regular maps of small genus}. The inspected material
does not independently supply that chiral census. We
retain the historical report with this explicit limitation.

The source's semisparseness criterion is likewise retained
with its actual quantifiers. Put
$\Gamma_{<i}=\gen{\rho_j:j<i}$,
$\Gamma_{>i}=\gen{\rho_j:j>i}$ and
$\Gamma_i=\gen{\rho_j:j\ne i}$.
For $\Sigma\le\Gamma$, the first two tests are
\begin{align*}
 \Sigma^\varphi\cap\Gamma_i\rho_i&=\varnothing,\\
 \Gamma_{>i}\cap\Gamma_{<j}\Sigma^\varphi
 &\subseteq\Gamma_k(\Gamma_{>i}\cap\Gamma_{<j})\Sigma^\varphi,
                    &&i<k<j,
\end{align*}
for every $\varphi\in\Gamma$ and all displayed valid indices.
The third is written in the source as
\[
 \Gamma_{>i}\cap\Gamma_{<i+1}\Sigma\subseteq\Sigma
                                      \quad(0\le i<n-1).
\]
Hartley's further necessary condition requires both end
intersections to be semisparse for every conjugate and
\[
 \Sigma^\varphi\cap\Gamma_0\Gamma_{n-1}
   =(\Sigma^\varphi\cap\Gamma_0)
      (\Sigma^\varphi\cap\Gamma_{n-1})
                   \quad(\varphi\in\Gamma).
\]
Problem~34 asks for sufficiency of these latter conditions.
The special locally projective census, with its additional
hypotheses, does not prove that general converse.

\section{The questions and the fixed edition}

The table below records every numbered problem. ``Unsettled
here'' states a limit of the present acquired proofs; it
does not assert that the question is still open in the
current literature. A broader search for later answers
is outside this intake. The fourteen secondary records
follow the table so that an embedded conjecture cannot
disappear within a broader programme.

\input{question_table.tex}
\input{secondary_table.tex}

Several edition qualifications matter for interpreting
these statements. Within the nondegenerate fundamental
domain used in the source, the chiral toroid condition is
$t\ne0$ \emph{and} $s\ne t$; the printed ``or'' would
include the regular cases. This correction alone does not
prove that every small parameter gives a nondegenerate
polytope. The reference to plane maps at rank $n=2$ is
inconsistent with the source's rank-three polyhedron
convention and is retained as an edition qualification.
Recursive face realization is used at all appropriate
ranks, rather than extrapolating the displayed $j=1,2,3$
notation without a definition. The abstract and
introduction give different meeting years; this historical
discrepancy does not change the fixed mathematical edition.
Three repeated TeX labels refer to different numbered
problems, so our register identifies problems by ordinal
number and source position, never by the duplicated label
alone.

\section{Proof record and attribution}

The three common theoretical components contain complete
proofs: four pages for finite realization cones, three
for rotation conventions and flag walks, and two for
conditional regular amalgamation. The two finite examples
add two pages with complete presentations, executable GAP
calculations, and independent literal permutation checks.
All four have distinct reviews. Their unabridged texts
are appended to the reading volume. The earlier group and
locally projective results are acquired with their full
proof and certificate dependencies, not just their final
orders or page counts.

The new work here is the arrangement of these proofs,
the explicit joins between their hypotheses, and the
short proof of Proposition~\ref{neighborly}. The source
authors and their cited predecessors retain attribution
for established statements. The source register separately
identifies definitions, retained classifications,
historical reports, acquired proofs, and questions.
No new open-problem solution is claimed. The computational
checks are not Lean certificates, and no new compression
factor is inferred from this survey disposition. Its
admission as a handled original requires the separate
review of this entire assembly.

\begin{thebibliography}{9}\small
\bibitem{survey} E. Schulte and A. Ivi\'c Weiss,
\emph{Problems on Polytopes, Their Groups, and Realizations},
Period. Math. Hungar. 53 (2006), 231--255;
arXiv:math/0608397v1. Fixed source and complete bibliography.
\bibitem{ARP} P. McMullen and E. Schulte,
\emph{Abstract Regular Polytopes}, Cambridge University
Press, 2002. Standard reconstruction and the classifications
at the locations specified in the retained-input register.
\bibitem{SW} E. Schulte and A. Ivi\'c Weiss,
\emph{Chirality and projective linear groups}, Discrete
Math. 131 (1994), 221--261. Retained through the
inspected survey, not independently reacquired here.
\bibitem{HW} I. Hubard and A. Ivi\'c Weiss,
\emph{Self-duality of chiral polytopes}, J. Combin.
Theory Ser. A 111 (2005), 128--136. Attribution and
oriented examples as cited in the fixed source.
\bibitem{components} Human Compression project,
\emph{Finite realization cones};
\emph{Rotation conventions and closed flag walks};
\emph{Universal regular amalgamation from a witness};
\emph{Two oriented square-torus amalgamations}.
Complete acquired manuscripts and distinct reviews.
\bibitem{acquired} The exact acquired group, chamber,
and locally projective manuscripts are listed with their
source attributions and immutable editions in the
acquired-proof register. Their original results receive
no new priority or question credit in this paper.
\end{thebibliography}
\end{document}
